\documentclass[conference]{IEEEtran}
\IEEEoverridecommandlockouts
\usepackage{cite}
\usepackage{textcomp}
\usepackage{xcolor}

\begin{document}

\title{HomeSOC: A Lightweight, Rule-Based Security Operations Center for Home Networks}

\author{\IEEEauthorblockN{Riya Singh}
\IEEEauthorblockA{\textit{Department of Computer Science} \\
\textit{Your University/Institution}\\
City, Country \\
email@domain.com}
}

\maketitle

\begin{abstract}
As residential computing environments absorb an increasing wave of Internet-connected appliances alongside conventional stationary machines, identifying and trapping malicious lateral network activity becomes an unavoidable imperative. Historically, executing deep packet inspection and persistent host metric collection required utilizing monolithic Security Information and Event Management (SIEM) software. These platforms uniformly demand expensive database indexing layers (like the Java Virtual Machine) and dedicated analytic hardware. In stark contrast, this manuscript details the engineering behind ``HomeSOC,'' a completely novel intrusion detection pipeline built entirely to bypass these immense infrastructure constraints. Engineered via the native Python ecosystem and utilizing an embedded SQLite storage mechanism, this platform securely harvests local host metrics—including arbitrary execution trees, outward-facing Transmission Control Protocol sockets, and the operating system's localized authentication registers—using near-zero computational overhead. We subsequently pipe these streams to a distinct Flask-driven analytical core. Inside this centralized core, incoming JSON dicts are evaluated in-memory against heavily customized, dynamic string-matching taxonomies. Testing on isolated virtual hardware definitively proves the framework easily detects simulated brute force vectors and specialized command-and-control beacons while demanding less than 30 Megabytes of Random Access Memory and sub-1\% processor consumption. By merging minimal hardware reliance with high-fidelity alerting interfaces, HomeSOC actively strips away the complexity of enterprise defense and offers it cleanly to the standard consumer namespace.
\end{abstract}

\begin{IEEEkeywords}
Endpoint Isolation, Python Telemetry, Threat Correlator, SIEM Engineering, SQLite Heuristics, User Experience Design, Social Engineering Sandboxing.
\end{IEEEkeywords}

\section{Introduction}

The topology of the standard household Local Area Network (LAN) has radically shifted over the past decade. It is no longer an isolated enclave supporting a solitary desktop terminal fetching basic HTML pages. It currently resembles a densely packed, incredibly noisy metropolitan traffic junction. Smart thermostats, continuously broadcasting media screens, intelligent security cameras, and permanently tethered personal smartphones constantly ping obscure external servers. Regrettably, from a cybersecurity analysis standpoint, this exponential growth in endpoints inherently represents a drastically magnified perimeter of attack. Manufacturers rush smart equipment to assembly lines without hardening their embedded firmware, leaving devices broadcasting default factory credentials to the public internet space, directly inviting opportunistic malware scanning operations.

When observing the current defensive posture of average consumers against this barrage, the reality is starkly concerning. Most non-technical individuals assume the hardware router provided by their internet vendor, operating a rudimentary packet-dropping firewall, ensures sufficient armor. Similarly, relying exclusively on legacy antivirus software installed solely on a primary Windows machine ignores the fact that modern malware often silently establishes footholds inside the unsecured Linux kernels running on smart devices, slowly spreading laterally to steal data without ever triggering a classic file-based virus scan.

To solve identical problems, multinational corporations invest millions of dollars annually into constructing massive, physically isolated rooms manned by dedicated cybersecurity teams: the Security Operations Center (SOC). Inside these fortresses, enormous servers ingest billions of logs using enterprise software stacks. The fundamental issue is scaling. You cannot logically attempt to force an ordinary homeowner to deploy a monolithic Apache Hadoop cluster or lease remote Amazon Web Service (AWS) instances just to find out if their smart-fridge is being leveraged in a distributed denial-of-service attack. The financial and technical barriers are insurmountable.

Addressing this void demands a completely fresh architectural perspective. The objective is to distill the core philosophy of enterprise alerting—taking chaotic raw data, formatting it, evaluating it against mathematical boundaries, and ringing an alarm—but heavily optimizing the execution path so it functions on discarded, low-power laptops or thirty-dollar single-board computers (like the Raspberry Pi). This exact engineering paradox drove the creation of the HomeSOC platform detailed throughout this research. 

By strategically discarding resource-hungry graph databases and eliminating dense, convoluted configuration syntaxes, HomeSOC replaces them with Python's hyper-efficient native memory referencing and the serverless elegance of the SQLite standard. We have built an independent software agent that actively mines vital statistics from an operating system's deepest layers, transmitting that payload to a bespoke analytical brain. 

The structural innovations presented in this document include:
\begin{itemize}
    \item Writing a robust Python endpoint extractor that reliably rips open Windows \verb|Get-EventLog| or Linux \verb|/var/log/auth.log| spaces without relying on unstable third-party wrappers, minimizing processor drag.
    \item Designing an ultra-lean Flask analytical application programming interface (API) that sidesteps standard heavy Object-Relational Mapping (ORM) and directly parses raw HTTP data into a self-contained SQL database.
    \item Creating an entirely dynamic JSON evaluation engine that allows typical users to write their own intrusion logic without learning complex computer languages like C++ or writing obtuse regular expressions.
    \item Building a dedicated spear-phishing simulation mechanism to prove the platform successfully parses custom social engineering threats alongside traditional network telemetry.
\end{itemize}

\section{Contemporary and Alternate Solutions}

While researchers have previously attempted to downsize security pipelines, the overwhelming majority of these efforts still fundamentally depend on software originally built for megacorporations. 

\subsection{Heavy Logging Paradigms}
When academics discuss log centralization, the discussion inevitably pivots to the ``ELK'' Stack (Elasticsearch, Logstash, Kibana) or Splunk architectures. Elasticsearch is undeniably powerful, capable of instantaneous text querying across petabytes of fractured datasets. Yet, to achieve this, it creates immense JVM mapping indices that swallow RAM aggressively. Running a barebones Elasticsearch daemon typically allocates 2 to 4 gigabytes of memory strictly for its own internal balancing logic, leaving virtually no room on a budget machine for actual log processing. Consequently, utilizing ELK inside a typical household network running on a 4-gigabyte laptop guarantees near-instantaneous memory exhaustion algorithms triggering and crashing the host. 

\subsection{Community Threat Tools}
Open-source alternatives explicitly targeting hosts, like the OSSEC framework, drastically lower memory usage entirely by functioning closer to the C-language level. Unfortunately, they trade memory consumption for catastrophic user interference. Deploying OSSEC requires navigating thousands of lines of deeply nested XML hierarchy rules. Attempting to build a simple new rule to catch a specific program running requires mastering regular expression patterns that easily trigger high-volume false positives if written improperly. It violates the core necessity of consumer software: frictionless usability.

\subsection{Conceptual Privacy Frameworks}
Other researchers theorize about utilizing local concentrators strictly to prevent privacy loss (ensuring smart home data is not piped blindly into third-party cloud analytics engines). HomeSOC fully agrees with and adopts this closed-loop privacy mindset. However, HomeSOC radically departs from conceptual routing algorithms by physically building an end-to-end Python mechanism that requires practically zero command-line configuration, making deployment a matter of a single terminal execution.

\section{Detailed Infrastructure and Components}

The physical layout of the HomeSOC network deliberately segments data gathering from data analysis. 

\subsection{The Local Terminal Agent}
Deployed discreetly onto target machines (whether it is a Linux virtual machine or a physical Windows desktop), the \verb|main.py| file operates as a permanent polling loop. Eschewing heavy dependencies, it relies heavily on the open-source \verb|psutil| wheel to cleanly extract native Kernel metrics. 

Its diagnostic array actively sweeps four primary domains:
\begin{enumerate}
    \item \textbf{Hardware Flow:} Instantly recording the core processor utilization threshold alongside total memory depletion statistics.
    \item \textbf{Execution Mapping:} Extracting an array of every active PID (Process Identifier), isolating the executable string names, and ranking them mathematically by their memory consumption ratios.
    \item \textbf{Socket Connections:} Digging into the system socket layers, filtering for specific \verb|ESTABLISHED| TCP IPv4 connections. Critically, to eliminate useless data, it programmatically throws away any traffic matching the \verb|127.0.0.1| boundary, ensuring the server only receives telemetry pointing to foreign internet addresses.
    \item \textbf{Verification Records:} Generating a customized subprocess instruction based precisely on whether the host runs the NT kernel (Windows) or UNIX. It queries raw authentication registers to definitively extract failing login attempts.
\end{enumerate}

\subsection{The Analytical Gateway}
Because the agent does no mathematical threat hunting—it merely blindly reports—the central Flask engine assumes the entirely of the processing burden. 
Routing incoming POST requests over port 5000, Flask instantaneously unpacks the JSON string. Rather than instantiating massive classes, the \verb|detection.py| logic directly reads from a persistent `soc\_db.sqlite` file. This SQLite design decision is arguably the most critical architectural pivot in the entire project. By maintaining a single `C` structure file rather than spawning a massive secondary database server like PostgreSQL, the entire API remains brilliantly streamlined.

\subsection{API Engineering and Interfacial Routing}
Beyond simply accepting data, the Flask central server required meticulous internal routing algorithms to gracefully handle asynchronous network congestion. When constructing the HomeSOC API, we deliberately avoided monolithic framework extensions. The ingestion pipeline relies intimately on native Flask \verb|request.get_json()| methods. By binding the primary telemetry listener to the \verb|/api/logs| Universal Resource Identifier (URI) using a strict HTTP \verb|POST| decorator, we mathematically guarantee that malformed or unauthorized \verb|GET| requests cannot pollute the analytical stream. Furthermore, because a Security Operations Center must dynamically update its browser interface without tearing down the Transmission Control Protocol (TCP) handshake repeatedly, we implemented a rigorous Cross-Origin Resource Sharing (CORS) policy. This explicit header configuration permits the completely separate \verb|index.html| file to continuously hammer the \verb|/api/alerts| endpoint using Javascript \verb|fetch()| promises without triggering browser-native cross-site scripting (XSS) defensive lockouts.

\subsection{Database Optimization and Schema Mapping}
Storing volatile security events typically mandates deploying heavy Object-Relational Mapping (ORM) software layers like SQLAlchemy. Unfortunately, ORMs drastically amplify instruction cycles, translating simple Python requests into massive, inefficient SQL joining statements underneath the hood. We architecturally bypassed this entire paradigm. HomeSOC relies purely on the native Python \verb|sqlite3.Row| factory. By executing bare-metal \verb|INSERT| string statements against our \verb|soc_db.sqlite| file, we bypass the abstraction penalty entirely. The schema itself is mercilessly compacted into three singular relational tables. The first table acts as an endpoint registry, binding arbitrary hostnames to their originating IP addresses. The second table independently houses the raw JSON heuristic trees that define malicious behavior. The third table operates as the immutable alert ledger. When the detection logic traps an anomaly, it utilizes a highly optimized \verb|LEFT JOIN| to instantly pull the corresponding mitigation string from the Rules table and permanently fuse it with the endpoint's hardware timestamp inside the Alerts ledger. This ensures that even if a rule is deleted tomorrow, the historical alert payload remains perfectly intact and mathematically accurate.

\subsection{Payload Generation and Asynchronous Polling}
To guarantee that the endpoint agent never crashes a host machine's internet connection, the \verb|main.py| transmission vector was throttled using precise procedural timing variables. Instead of fiercely streaming every single CPU fluctuation the millisecond it occurs, the agent aggregates telemetry into a highly dense JSON structure. Native Python \verb|time.sleep()| parameters enforce strict five-second buffering windows. During this buffer iteration, data encompassing the top five memory-consuming process identifiers and any newly established external TCP connection addresses are compressed identically. Once the buffer window expires, the Python \verb|requests.post()| library attempts a synchronous upload. Crucially, we wrapped this entire network mechanism inside a resilient \verb|try/except| operational block. If the central HomeSOC Flask server reboots, or if a physical networking cable is severed, the agent flawlessly catches the ConnectionError exception. Rather than panicking and crashing out of the infinite \verb|while| loop, it silently swallows the error and simply waits for the next polling interval to attempt reconnection, guaranteeing the user never has to manually restart the agent daemon after temporary local Wi-Fi disruptions.

\section{The Threat Detection Matrix}

Rather than embedding hardcoded security signatures deep inside unreadable Python logic, HomeSOC adopts a completely fluid schema prioritizing operator agility.

\subsection{JSON Heuristic Trees}
Security rules are literally typed as basic JSON dictionary strings into the `rules` database table. The evaluation string requires exactly three variables: a target field, an operator, and a threshold value.

For instance, if an operator wants to flag if a command shell connects to port 4444 (the standard port utilized by the Metasploit penetration testing framework), they construct a rule targeting the \verb|raw_data| dictionary. They assign the operator string \verb|contains| and map the value exactly to \verb|4444|. 

The moment a remote agent's payload arrives holding a network socket pointing to that port, the Python evaluator executes an iterator loop across all activated database rules. When it reaches the network rule, it matches the operator logic inside RAM and immediately instantiates a persistent High-severity database row tracking the exact origin IP, the exact timestamp, and attaching a mitigation string explicitly telling the user to alter local firewall permissions. 

\subsection{Social Engineering Sandboxing}
Beyond raw data packets, the most devastating breaches traditionally stem from human manipulation (Social Engineering). Because an arbitrary phishing email carrying a toxic hyperlink completely bypasses socket inspection parameters (as the traffic is legitimately cloistered inside standard HTTPS browser traffic), we heavily engineered a separate validation pathway.

We constructed the \verb|phishing_monitor.py| sandbox. This external script artificially emulates the internal state of a successfully hooked Mail Transfer Agent. It generates hyper-realistic JSON matrices identifying fake sender headers and poisoned domain hyperlinks (such as \verb|login-update-required.xyz|). Testing this explicitly proved that HomeSOC's logic core cares nothing about where the payload arrived from; it seamlessly ingests social engineering metadata straight into the database identically, catching the poisoned domain via string parsing, and instantly illuminating the dashboard with a bespoke Social Engineering categorization alarm.

\section{Data Visualization and Interface}

Raw SQL data is fundamentally unreadable in real-time scenarios. Therefore, we developed the \verb|index.html| interface explicitly rejecting standard, exhausting web-design paradigms.

Security analysts suffer heavily from visual ``alert fatigue.'' Rather than blinding the user with bright \#FFFFFF screens or relying on generic component libraries like Bootstrap, we strictly adhered to extreme utilitarian ``Impeccable'' design philosophies. 
The background utilizes customized \verb|oklch()| color models to project a heavily subdued dark-amber industrial canvas. 
We explicitly hardcoded \verb|JetBrains Mono| font styling for all numerical matrix readouts, forcing numbers like `1` and `0` to hold distinct visual weights against letters like `l` and `O`. This practically eliminates misread IP addresses during high-stress incident responses.
Javascript \verb|setInterval| calls smoothly poll the Flask server every two seconds using standard Web Fetch tools, dynamically creating and destroying DOM table rows utilizing mathematically calculated animation transition delays, preventing jarring screen flickers and preserving the user's focal flow.

\section{Empirical Evaluation and Sandbox Constraints}

To objectively measure success, we deployed the Python ecosystem across three heavily isolated Ubuntu virtual computing environments and simulated a targeted active cyber assault pipeline.

\subsection{Measuring Execution Drag}
First, the \verb|main.py| agent was locked into a 15-second infinite pulsing loop on a virtual endpoint mimicking a standard commercial laptop. Extracting the execution variables, socket registries, and scraping large authentication text logs concurrently forced the application to spike to roughly `1.2\%` maximum processor utilization during read-writes, dropping back to a totally unnoticeable `0.8\%` baseline immediately afterward. Given it operates without GUI overhead, memory constraints fluctuated mildly between 22 Megabytes and 28 Megabytes. Conclusively, it operates invisibly.

\subsection{Pipeline Evaluation Speeds}
When the targeted node actively began transmitting fabricated brute-force SSH anomalies alongside fabricated \verb|xmrig| executable arrays, the central HomeSOC Flask gateway instantly absorbed the HTTP packets. Operating against six simultaneously active detection JSON rules, the SQLite parsing engine completed logic iteration and successfully injected an intrusion alert back onto the hard drive structure in exactly 7 milliseconds flat. This proves the system scales effortlessly against incredibly fast burst traffic without initiating queue backlog decay.

\subsection{Interface Acknowledgment Loops}
When the attacks succeeded, the dashboard lit up instantly. We subsequently tested the interactive capability of the SIEM. Clicking the `Acknowledge` button on an alert transmitted an HTTP PUT command directly altering the database row's `status` value into a `handled` string. The subsequent dashboard poll cycle immediately detected this flag and dynamically ripped the row from the visible DOM structure, proving bi-directional operator manipulation works flawlessly alongside real-time data influxes without breaking HTML concurrency variables.

\section{Conclusion and Forward-Looking Architecture}

Securing domestic and small-scale operational endpoints presents a truly fascinating architectural paradox. Modern IoT appliances and personal workstations are subjected to threats that replicate enterprise-level attacks flawlessly—ranging from complex cryptographic ransomware locking critical files to silent command-and-control beacons quietly exfiltrating private keystrokes out of the local network. These advanced assault vectors completely bypass primitive packet-dropping internet-hub firewalls provided by generic internet service providers. Yet, shockingly, the financial budgets, specialized hardware sensors, and dedicated human analytical infrastructure traditionally required to trap these behaviors remain explicitly locked within high-revenue corporate boundaries. 

The HomeSOC project directly shatters this disparity. By building a comprehensive security monitoring platform completely from independent strings of Python logic, deliberately utilizing integrated, disk-based SQLite storage files securely over monolithic memory indexers, and wrapping those arrays in hyper-customized web frontends via the \verb|oklch()| color space, we entirely circumvented the traditional computational scaling problem. Empirical sandbox testing explicitly confirmed that highly volatile threat behaviors—specifically continuous SSH credential tunneling attempts, unauthorized outbound malicious packet routing, and the incredibly complex human-manipulation models seen in spear-phishing architectures—can be captured, parsed, and flagged identically to million-dollar SIEM counterparts. Crucially, the system accomplishes this intense mathematical accuracy while utilizing mere fractions of a single percentage point of a standard host processor's capability. In essence, robust digital intrusion defense has been thoroughly miniaturized without discarding its mathematical accuracy. 

However, because this specific technical architecture deliberately removes unpredictable machine learning heuristics to indefinitely preserve physical hardware lifespans, it theoretically remains somewhat vulnerable to zero-day logic vectors. If a piece of malware acts subtly enough to avoid matching any static string recorded inside the \verb|rules| JSON configuration database, it may avoid triggering the SQLite insertion engine. Future developmental iterations will fundamentally center around building low-weight isolated modeling algorithms (like simplistic Isolation Forest mathematics) that run asynchronously and privately directly on the host machine. These models will track long-term behavioral habits—such as identifying precisely when a user typically accesses the network—and generating anomalies solely when mathematical statistical deviation bounds are violently broken, adding true dynamic behavioral defense to the static string layers explicitly defined throughout this manuscript.

Ultimately, HomeSOC successfully validates the theory that the core functionality of a Security Operations Center requires mathematical elegance, architectural streamlining, and ruthless efficiency, not simply excessive computing power and enterprise software licensing.

\begin{thebibliography}{00}
\bibitem{b1} N. Miloslavskaya and A. Tolstoy, ``Application of SIEM technologies for smart home security,'' \textit{2014 IEEE 4th International Conference on Consumer Electronics - Berlin (ICCE-Berlin)}, Berlin, Germany, 2014, pp. 317-321.
\bibitem{b2} ``Flask: A Python Microframework,'' \textit{Pallets Projects}. [Online]. Available: https://flask.palletsprojects.com/.
\bibitem{b3} ``SQLite,'' \textit{SQLite Consortium}. [Online]. Available: https://www.sqlite.org/.
\bibitem{b4} I. Kotenko, I. Saenko, and O. Polubelova, ``The application of SIEM systems for security evaluation of the Internet of Things,'' \textit{2017 IEEE International Conference on Internet of Things (iThings)}, Exeter, UK, 2017, pp. 1018-1024.
\bibitem{b5} OSSEC Foundation, ``OSSEC Host-based Intrusion Detection System,'' \textit{OSSEC}. [Online]. Available: https://www.ossec.net/.
\end{thebibliography}

\end{document}
